\chapter*{Resumen}
\thispagestyle{empty}

La presente tesis describe el desarrollo de un prototipo
de sistema para el reconocimiento automático de números de placa (ANPR),
pensado para su futura integración en un sistema de gestión inteligente de estacionamientos
en el Campus Tecnológico Local San José.

El diseño del sistema ANPR plantea en dos etapas:
detección de la placa vehicular y reconocimiento de caracteres.
Para la detección se emplea la arquitectura YOLO, 
ajustada para detectar la región en la que se ubica la placa de los vehículos.
Para el reconocimiento se utiliza Fast-Plate-OCR,
adaptado al estilo tipográfico y variedad de placas costarricenses.
Dado el tamaño limitado del conjunto de datos, se desarrolla un generador
de imágenes sintéticas y un esquema de aumento de datos 
para entrenar el módulo de reconocimiento de caracteres.
Finalmente, se realizaron experimentos para evaluar del desempeño tanto de cada componente del
sistema por separado como del sistema integrado.

El sistema desarrollado alcanzó $98.4\%$ de precisión en el 
reconocimiento en placas vehiculares costarricenses
y $93.5\%$ de precisión de la detección de la región donde 
se encuentra la placa vehicular.
El sistema integrado procesa video a $3.84\pm 0.79$ FPS,
con una latencia media de $299\pm 69.7$ ms.

\bigskip

%% Defina las palabras clave con defKeywords en config.tex:
\textbf{Palabras clave:} ANPR, IPLMS, extremo a extremo, imagen sintética \thesisKeywords

\clearpage
\chapter*{Abstract}
\thispagestyle{empty}

This thesis describes the development of a prototype system for automatic license plate recognition (ANPR),
designed for future integration into an intelligent parking management system 
at the San José Local Technology Campus.

The ANPR system design consists of two stages: vehicle license plate detection and character recognition.
The YOLO architecture is used for detection, adjusted to detect the region where the vehicle license plate is located.
Fast-Plate-OCR is used for recognition, adapted to the typographic style and variety of Costa Rican license plates.
Given the limited size of the dataset, a synthetic image generator and a data augmentation scheme 
were developed to train the character recognition module.
Finally, experiments were conducted to evaluate the performance of each component of the system separately 
and of the integrated system.

The developed system achieved 98.4\% accuracy in recognizing Costa Rican license plates
and 93.5\% accuracy in detecting the region where the license plate is located.
The integrated system processes video at $3.84\pm 0.79$ FPS, with an average latency of $299\pm 69.7$ ms.

\bigskip

\textbf{Keywords:} ANPR, IPLMS, end-to-end, synthetic images

\cleardoublepage

%%% Local Variables: 
%%% mode: latex
%%% TeX-master: "main"
%%% End: 
